\section{ICE, STUN, TURN và topology media}

\subsection{NAT traversal}

ICE tìm đường truyền qua NAT bằng các candidate \cite{rfc8445}. Host candidate là địa
chỉ cục bộ. Server-reflexive candidate là địa chỉ công khai do STUN phát hiện. Relay
candidate là địa chỉ chuyển tiếp do TURN cấp. Hai phía trao đổi và kiểm tra các candidate
qua signaling.

TURN chuyển tiếp media khi đường truyền trực tiếp thất bại \cite{rfc8656}. Môi trường
production yêu cầu TURN do hệ thống kiểm soát, credential có thời hạn, TLS và giám sát
băng thông.

\begin{figure}[htbp]
  \centering
  \includegraphics[width=\textwidth]{webrtc-arch-custom.png}
  \caption{Signaling, STUN/TURN và các đường truyền media trong WebRTC.}
  \label{fig:webrtc-architecture}
\end{figure}

Kiểm thử TURN phải đặt \texttt{iceTransportPolicy=relay} hoặc xác nhận selected candidate
pair có loại \texttt{relay}. WebRTC Statistics API cung cấp candidate pair, RTT, jitter,
packet loss, bitrate và thông tin frame \cite{w3c-stats}.

\subsection{P2P và SFU}

Trong mô hình P2P, hai trình duyệt truyền media trực tiếp. TURN tiếp nhận media khi P2P
thất bại. Ứng dụng tự quản lý khôi phục kết nối, quan sát và tương thích trình duyệt.

Trong mô hình SFU, mỗi client gửi media tới máy chủ và nhận track do máy chủ chuyển tiếp.
SFU tập trung việc điều phối track, quan sát, mở rộng và recording. Đổi lại, máy chủ phải
xử lý băng thông media. Với cuộc gọi 1:1, quyết định dùng SFU dựa trên yêu cầu vận hành
và mở rộng.
